\documentclass[11pt,a4paper]{article}
\usepackage[T1]{fontenc}\usepackage[utf8]{inputenc}
\usepackage[a4paper,top=2.6cm,bottom=2.6cm,left=2.7cm,right=2.7cm]{geometry}
\usepackage{graphicx}\usepackage{booktabs}\usepackage{array}\usepackage{enumitem}
\usepackage{caption}\usepackage{fancyhdr}\usepackage{amsmath}\usepackage{parskip}
\usepackage[hidelinks]{hyperref}\usepackage{float}
\captionsetup{font=small,labelfont=bf,labelsep=period,justification=raggedright,singlelinecheck=false}
\setlist[itemize]{leftmargin=1.3em,itemsep=4pt,topsep=4pt}
\setlist[enumerate]{leftmargin=1.6em,itemsep=6pt,topsep=4pt}
\pagestyle{fancy}\fancyhf{}
\fancyhead[L]{\footnotesize KOUAME KOFFI FID\`ELE}
\fancyhead[C]{\footnotesize TASK 2 \textendash{} SALES ANALYSIS}
\fancyhead[R]{\footnotesize Data Analysis Internship}
\fancyfoot[C]{\thepage}\renewcommand{\headrulewidth}{0.4pt}
\begin{document}

\begin{titlepage}\centering
\IfFileExists{logos.png}{\includegraphics[width=\linewidth]{logos.png}\\[1.4cm]}{\vspace*{0.4cm}}
\rule{\linewidth}{0.5pt}\\[0.9cm]
{\large\bfseries Data Analysis Internship}\\[0.7cm]
{\normalsize\scshape Task 2 Report}\\[1.5cm]
{\LARGE\bfseries Sales Data Cleaning, Validation}\\[0.55cm]
{\LARGE\bfseries and Performance Dashboard}\\[1.6cm]
\rule{\linewidth}{0.5pt}\\[1.5cm]
{\large\itshape Analysis carried out on}\\[0.45cm]
{\Large\bfseries THE FACTSALE DATASET}\\[0.45cm]
{\normalsize 26{,}397 order lines, January 2013 to May 2016}\\[1.5cm]
{\large\itshape Prepared by}\\[0.45cm]
{\Large\bfseries KOUAME KOFFI FID\`ELE}\\[1.6cm]
{\large\bfseries \today}\vfill\end{titlepage}

\tableofcontents\newpage

\section{Executive summary}

The brief anticipates a dirty file: impossible calendar dates, duplicate rows, broken
arithmetic. Each of those was tested explicitly. \textbf{Most of them are not present.}
Every arithmetic identity in the file holds to within one cent, there are no duplicate
records, and every date parses without exception.

Reporting corrections that were not needed would have been the easy path. This report
does the opposite: it states what the file actually contains, and the three findings
that matter are structural rather than arithmetic.

\begin{enumerate}
\item \textbf{9{,}077 order lines, 34.4\% of the file, carry no customer key.}
Every customer metric therefore describes two thirds of the business.
\item \textbf{The last year is incomplete}, January to May 2016 only. Any year-on-year
chart built without noticing shows a collapse that never happened.
\item \textbf{A routine outlier screen would have deleted 48.0\% of all revenue.}
Sales value is right-skewed by construction, and a symmetric rule applied to a skewed
variable removes the top of the distribution rather than errors.
\end{enumerate}

No row was deleted. All 26{,}397 lines are retained.

\section{Key performance indicators}

\begin{table}[H]\centering
\caption{The seven indicators selected for the dashboard}\small
\begin{tabular}{@{}lrp{6.4cm}@{}}\toprule
\textbf{KPI} & \textbf{Value} & \textbf{Why it was selected} \\ \midrule
Total sales, excl. tax & 19.88M & The top line every other figure divides into \\
Gross profit & 9.92M & Revenue that does not convert is not performance \\
Gross margin & 49.9\% & The one ratio that survives a change of scale \\
Orders (invoices) & 8{,}188 & Separates growth in volume from growth in basket \\
Average order value & 2{,}428 & The lever a commercial team can actually pull \\
Units sold & 1{,}028{,}670 & Detects margin bought with discount \\
Identified customers & 48 & Bounds every customer conclusion in this report \\
\bottomrule\end{tabular}\end{table}

\section{The dashboard}

\begin{figure}[H]\centering
\includegraphics[width=\linewidth]{07_dashboard.png}
\caption{The sales performance dashboard. An interactive version, with hover detail,
zoom and image export, is supplied as \texttt{dashboard.html} and opens in any browser
without a server or an installation.}
\end{figure}

\subsection{Why each chart type was chosen}

\begin{table}[H]\centering\small
\begin{tabular}{@{}p{3.6cm}p{3.0cm}p{6.6cm}@{}}\toprule
\textbf{Panel} & \textbf{Type} & \textbf{Reason} \\ \midrule
Sales trend & Line with moving average & Continuous time axis; the twelve-month window spans exactly one seasonal cycle, so the smoothed series isolates the trend \\
Seasonality & Vertical bars & Months are ordered categories; this answers which month sells most, which the time series cannot because each month appears several times \\
Package & Horizontal bars, log scale & One package type carries 96.8\% of sales; on a linear axis the other three are invisible \\
Top products & Horizontal bars & Long product names read horizontally; only the top ten, since 227 bars answer nothing \\
Customer concentration & Cumulative curve & Concentration is a distributional property, which a ranking cannot show \\
\bottomrule\end{tabular}\end{table}

A pie chart was deliberately not used. Comparing angles is less accurate than comparing
lengths, and it fails entirely on a distribution as skewed as the package split.

\section{Data quality}

\subsection{Checks that passed}

\begin{table}[H]\centering\small
\begin{tabular}{@{}lr@{}}\toprule
\textbf{Check} & \textbf{Result} \\ \midrule
Duplicate rows, full record & 0 \\
Duplicate business keys & 0 \\
Unparseable invoice dates & 0 \\
Quantity $\times$ Unit Price $\neq$ Total excl. tax & 0 \\
Total excl. $\times$ rate $\neq$ Tax amount & 0 \\
Total excl. $+$ Tax $\neq$ Total incl. & 0 \\
Quantity $\leq$ 0 & 0 \\
Negative sales & 0 \\
Delivery before invoice & 0 \\
\bottomrule\end{tabular}\end{table}

\subsection{Issues found and how each was handled}

\paragraph{Unidentified customers.} 9{,}077 lines carry a customer key of zero,
which is a placeholder rather than a customer. The rows are \textbf{kept} for all
revenue, product and time analysis, because their financial fields are complete and
removing them would delete a third of the revenue. They are \textbf{excluded} from
customer analysis only, through an explicit flag rather than a silent filter.

\paragraph{Incomplete final year.} The file ends on 31 May 2016. Seasonality and every
month-on-month comparison are computed on the complete years only, and the partial
period is shaded in every time series so a reader cannot mistake it.

\paragraph{Missing delivery dates.} 13 lines have none. They are
\textbf{not imputed}: the field feeds a delivery-lag metric, and an invented date would
manufacture a lag that never happened. The rows are retained, since every revenue field
on them is intact.

\paragraph{Loss-making lines.} 566 lines return a negative margin.
These are \textbf{kept}. A negative margin on a real sale is a commercial event, not a
data error, and removing it would inflate total profit.

\paragraph{High unit prices.} 113 lines exceed a unit price of 1{,}000.
Inspection shows every one is the same product, an air cushion machine at a consistent
price, with plausible quantities and sound arithmetic. It is an expensive product, not a
typing error, and it is \textbf{kept}.

\subsection{The outlier trap}

An interquartile screen at 1.5 times the range flags 2{,}337 lines on sales
value, that is 8.9\% of the file, and those lines carry
\textbf{48.0\% of all revenue}.

The screen was computed and \textbf{deliberately not applied}. The rule assumes an
approximately symmetric distribution. Sales value is strongly right-skewed by
construction, because a few large orders always exist. Applying a symmetric rule to a
skewed variable does not remove errors; it removes the top of the distribution, which is
precisely the revenue a sales analysis exists to explain.

\section{Insights and recommendations}

\subsection{A third of revenue cannot be attributed to a customer}

34.4\% of order lines carry no customer key. This is a data capture
failure upstream, and it is the most valuable finding of this audit: no segmentation, no
lifetime value and no retention measure can be trusted until it is closed.

\textbf{Recommendation.} Trace the unattributed lines back to their source system and
establish why the key is absent. Until then, publish every customer figure with its
coverage rate beside it.

\subsection{Products are concentrated, customers are not}

The top ten products carry 32.8\% of revenue out of 227 items.
Among the 48 identified customers, however, the top 20\% carry only
25\% of identified revenue, against 20\% under perfect equality.

Revenue is spread almost evenly across the account base, which is the opposite of the
pattern most retail businesses show. There is no key-account dependency, so no single
customer loss threatens the top line; but there is equally no high-value tier to build a
retention programme around.

\textbf{Recommendation.} Do not build a key-account programme on this base, because the
data says there are no key accounts. Concentrate commercial effort on the product side,
and protect the availability and margin of the top ten items before anything else.

\subsection{Price and volume move against each other, weakly}

The correlation between unit price and quantity is -0.14, and between
quantity and line revenue 0.24. The negative sign is expected:
expensive items sell in smaller quantities. Its weakness matters more. Revenue per line is
driven by price far more than by volume.

\textbf{Recommendation.} Do not use volume discounting as a revenue lever on this
catalogue. The correlation structure says it will not recover the revenue it gives up.

\subsection{Seasonality is real but moderate}

On the complete years the strongest calendar month outsells the weakest by a factor of
1.33. That is worth planning around, but it does not justify large
seasonal stock swings.

\textbf{Recommendation.} Align staffing and inventory to the quarterly pattern rather
than the monthly one, whose amplitude does not warrant the operational cost of adjusting
to it.

\subsection{Loss-making lines are a control issue, not a strategic one}

566 lines return a negative margin, costing 40{,}502.
Against total profit that is 0.41\%. Reporting it as a headline
would be an error of proportion.

\textbf{Recommendation.} Add an automated flag on negative-margin lines at invoicing
rather than opening a pricing review. The cost of the fix should match the size of the
problem.

\section{Limitations}

The dataset contains no cost of acquisition and no marketing spend, so customer
acquisition cost and lifetime value cannot be computed and are not reported. It contains
no geographic dimension beyond a city key that is not resolved to a place name, so no
regional analysis is offered.

The figures are descriptive. No statistical test accompanies them. The differences
discussed are large enough that this is unlikely to change the conclusions, but the point
is stated rather than assumed.

\section*{Deliverables}
\addcontentsline{toc}{section}{Deliverables}

\begin{table}[H]\centering\small
\begin{tabular}{@{}ll@{}}\toprule
\textbf{File} & \textbf{Content} \\ \midrule
\texttt{dashboard.html} & Interactive dashboard, opens in any browser \\
\texttt{report.pdf} & The present report \\
\texttt{Task2\_Analysis.ipynb} & Commented analysis notebook \\
\texttt{DATA\_QUALITY.md} & Full quality audit \\
\texttt{INSIGHTS.md} & Insights and recommendations \\
\texttt{data/FactSale\_clean.csv} & Cleaned dataset \\
\texttt{charts/} & Seven figures, PNG \\
\texttt{results/} & KPI and aggregate tables, plus the audit JSON \\
\bottomrule\end{tabular}\end{table}

\vspace{1.2cm}\noindent\rule{\linewidth}{0.4pt}\\[3pt]
{\small\textbf{KOUAME Koffi Fid\`ele}\\ Data Analysis Internship\\ koffifidelek59@gmail.com}
\end{document}
